\documentclass[11pt, a4paper]{article}

% --- Packages ---
\usepackage{tabularx}
\usepackage{array}
\usepackage{tikz}
\usepackage{pgfplots}
\pgfplotsset{compat=1.18} % Mengunci versi pgfplots agar stabil
\usetikzlibrary{shapes.geometric, arrows.meta, positioning}
\usepackage[utf8]{inputenc}
\usepackage[indonesian, english]{babel}
\usepackage[margin=1.8cm]{geometry} % Margin sedikit dirapatkan dari 2cm ke 1.8cm
\usepackage{amsmath, amsfonts, amssymb}
\usepackage{graphicx}
\usepackage{booktabs}
\usepackage{cite}
\usepackage{titlesec}
\usepackage{abstract}
\usepackage{hyperref}

% --- Setup Margin Abstrak Sama Presisi dengan Margin Halaman ---
\setlength{\absleftindent}{0pt}
\setlength{\absrightindent}{0pt}
\addto\captionsindonesian{\renewcommand{\abstractname}{Abstrak}}

\hypersetup{
    colorlinks=true,
    linkcolor=blue,
    citecolor=blue,
    urlcolor=blue
}

\titleformat{\section}{\large\bfseries}{\thesection.}{0.5em}{}
\titleformat{\subsection}{\normalfont\bfseries}{\thesubsection}{0.5em}{}

\setlength{\columnsep}{0.6cm}

\begin{document}

% --- Single Column (Full Page) ---
\begin{center}
      % --- Indonesian Title & Authors ---
      {\Large \bfseries Penguatan Ketangguhan NIDS dengan Reduksi Fitur Berbasis XGBoost Gain-based Importance melalui Adversarial Training terhadap Serangan Evasion Saliency Map \par}
      \vspace{0.2cm}
      {\small Nama Penulis$^{1}$, Penulis Kedua$^{2}$ \par}
      {\small $^{1,2}$Program Studi Teknik Informatika, Politeknik Elektronika Negeri Surabaya \par}
      {\small Email: \texttt{author@email.com} \par}
      \vspace{0.3cm}
\end{center}

  % --- Indonesian Abstract & Keywords ---
  \selectlanguage{indonesian}
  \begin{abstract}
  \small
  \noindent Keberhasilan Machine Learning dalam Network Intrusion Detection System (NIDS) seringkali terbatas pada kondisi trafik laboratorium yang bersih dan rentan terhadap manipulasi sengaja pada fase pengujian. Di sisi lain, optimasi model melalui reduksi dimensi fitur diperlukan untuk efisiensi sistem, namun berpotensi memperluas permukaan serangan (\textit{attack surface}). Penelitian ini bertujuan untuk mengevaluasi dan meningkatkan ketangguhan (\textit{robustness}) model XGBoost yang telah dioptimasi fiturnya menggunakan metode \textit{Gain-based Importance}. Penelitian ini menggunakan dataset CSE-CIC-IDS2018 yang mencakup 15,7 juta sampel trafik jaringan dari tujuh kategori serangan. Metodologi yang diusulkan mencakup tiga tahap utama: pertama, reduksi fitur dari 78 menjadi 10 fitur menggunakan metrik \textit{Gain} XGBoost; kedua, simulasi serangan \textit{evasion} berbasis \textit{Saliency Map} dengan variasi epsilon ($\epsilon \in \{0.01, 0.05, 0.1\}$); dan ketiga, implementasi \textit{Adversarial Training} sebagai mekanisme pertahanan. Hasil eksperimen menunjukkan bahwa model standar mengalami penurunan drastis pada \textit{Matthews Correlation Coefficient} (MCC) dari 0,9351 menjadi 0,0184 (penurunan 98\%) saat menghadapi serangan \textit{evasion} ($\epsilon = 0{,}1$). Namun, model yang diperkuat melalui \textit{Adversarial Training} berhasil memulihkan MCC hingga 0,9953 pada data adversarial sambil mempertahankan integritas deteksi pada trafik normal (MCC = 0,9347, penurunan hanya 0,04\% dari baseline). Analisis \textit{Robustness Ablation Study} pada empat konfigurasi fitur (Full/68, Top-15, Top-10, Top-5) mengkonfirmasi bahwa konfigurasi Top-10 merupakan titik optimal (\textit{sweet spot}) yang memberikan keseimbangan terbaik antara efisiensi komputasi dan ketangguhan terhadap serangan \textit{evasion}.
  \end{abstract}

  \vspace{0.15cm}
  \noindent\small\textbf{Kata Kunci:} NIDS, XGBoost Gain-based Importance, Reduksi Fitur, Adversarial Training, Evasion Attack, Saliency Map, CSE-CIC-IDS2018.

  \vspace{0.35cm}

  % --- English Title, Abstract & Keywords ---
  \selectlanguage{english}
  \begin{center}
      {\Large \bfseries Strengthening NIDS Resilience with XGBoost Gain-based Importance Feature Reduction via Adversarial Training against Saliency Map Evasion Attacks \par}
      \vspace{0.2cm}
  \end{center}

  \begin{abstract}
  \small
  \noindent The success of Machine Learning in Network Intrusion Detection Systems (NIDS) is often limited to clean laboratory traffic conditions and remains vulnerable to intentional manipulation during the testing phase. On the other hand, model optimization through feature dimension reduction is required for system efficiency; however, it has the potential to expand the attack surface. This study aims to evaluate and enhance the robustness of an XGBoost model whose features have been optimized using the Gain-based Importance method. The study utilizes the CSE-CIC-IDS2018 dataset comprising 15.7 million network traffic samples spanning seven attack categories. The proposed methodology comprises three main stages: first, feature reduction from 78 to 10 features using the XGBoost Gain metric; second, simulation of Saliency Map-based evasion attacks with epsilon variations ($\epsilon \in \{0.01, 0.05, 0.1\}$); and third, the implementation of Adversarial Training as a defense mechanism. Experimental results demonstrate that the baseline model suffers a drastic Matthews Correlation Coefficient (MCC) drop from 0.9351 to 0.0184 (a 98\% decline) when subjected to evasion attacks ($\epsilon = 0.1$). However, the model hardened through Adversarial Training recovers MCC to 0.9953 on adversarial data while maintaining detection integrity on normal traffic (MCC = 0.9347, a loss of only 0.04\% from baseline). The Robustness Ablation Study across four feature configurations (Full/68, Top-15, Top-10, Top-5) confirms that the Top-10 configuration represents the optimal sweet spot, providing the best balance between computational efficiency and resilience against evasion attacks.
  \end{abstract}

  \vspace{0.15cm}
  \noindent\small\textbf{Keywords:} NIDS, XGBoost Gain-based Importance, Feature Reduction, Adversarial Training, Evasion Attack, Saliency Map, CSE-CIC-IDS2018.

  \vspace{0.4cm}

% --- Main Content ---
\selectlanguage{indonesian}

\section{Pendahuluan}


\section{Metode Penelitian}
Metodologi penelitian ini dirancang secara sistematis untuk mengevaluasi dan meningkatkan ketangguhan (\textit{robustness}) \textit{Network Intrusion Detection System} (NIDS) terhadap serangan \textit{evasion}. Secara garis besar, metodologi ini mencakup enam tahapan utama seperti yang diilustrasikan pada Gambar \ref{fig:research_flowchart} dan dijelaskan pada sub-bagian berikut.

% --- FLOWCHART KESELURUHAN METODOLOGI ---
\begin{figure}[htbp]
\centering
\includegraphics[width=0.45\linewidth]{alur2.png}
\caption{Diagram Alur Metodologi Penelitian}
\label{fig:research_flowchart}
\end{figure}

\subsection{Dataset dan Pra-pemrosesan Data}
Penelitian ini menggunakan dataset CSE-CIC-IDS2018 yang disediakan oleh \textit{Canadian Institute for Cybersecurity} (CIC). Dataset ini dipilih karena memiliki cakupan serangan yang sangat luas dan mencerminkan skenario jaringan modern, yang terdiri dari 10 file CSV dengan total ukuran data mencapai kurang lebih 6,7 GB. Dataset ini mencakup profil trafik normal (\textit{benign}) serta tujuh kategori serangan utama, yaitu: \textit{Brute-Force} (SSH/FTP), DoS, DDoS, \textit{Web Attacks}, \textit{Infiltration}, dan \textit{Botnet}.

Mengingat volume data yang sangat masif, khususnya pada file trafik tanggal 20 Februari yang mencapai 3,8 GB, penelitian ini menerapkan strategi \textit{Stratified Sampling} sebesar 10\% untuk menjaga distribusi label serangan asli tetap representatif dalam subset data yang digunakan. Hasil sampling menghasilkan 1.570.577 sampel dari total 15.705.786 sampel populasi asli. Proses pembersihan data (\textit{data cleaning}) dilakukan secara sistematis melalui beberapa tahapan krusial untuk menjamin integritas model:

\begin{enumerate}
    \item \textbf{Penanganan Nilai Singular:} Semua sampel yang mengandung nilai \textit{Missing Values} (NaN) dihapus secara permanen dari dataset. Nilai \textit{Infinity} pada fitur numerik ditangani dengan metode imputasi menggunakan nilai finit maksimal pada kolom terkait guna mencegah kegagalan matematis pada saat pelatihan model.
    \item \textbf{Eliminasi Fitur Bias:} Fitur-fitur identitas aliran (\textit{flow identifiers}) seperti \textit{Flow ID}, \textit{Source IP}, \textit{Destination IP}, \textit{Source Port}, \textit{Destination Port}, dan \textit{Protocol} dibuang karena tidak memberikan kontribusi pada generalisasi pola serangan dan berisiko menimbulkan bias pelatihan. Kolom \textit{Timestamp} juga dihilangkan untuk memastikan model tidak mempelajari urutan waktu melainkan karakteristik statistik trafik.
    \item \textbf{Reduksi Variansi Rendah:} Dilakukan pemindaian terhadap fitur-fitur dengan \textit{zero-variance} (nilai konstan), di mana fitur yang tidak memiliki variabilitas data dibuang karena tidak mengandung informasi diskriminatif bagi algoritma klasifikasi.
    \item \textbf{Normalisasi dan Encoding:} Label serangan yang bersifat kualitatif dikonversi menjadi representasi numerik menggunakan \texttt{LabelEncoder}. Fitur numerik lainnya dinormalisasi menggunakan \texttt{StandardScaler} atau \texttt{MinMax Scaling} untuk memastikan distribusi data berada pada rentang yang seragam, sehingga mempercepat konvergensi algoritma XGBoost.
\end{enumerate}

Hasil akhir dari proses ini adalah file yang telah dibersihkan (\texttt{cleaned\_100.pkl/csv}) yang mengandung 68 fitur trafik teknis yang siap digunakan sebagai basis untuk tahap seleksi fitur berbasis \textit{Gain-based Importance}.


% Bab 2.2
\subsection{Optimasi Fitur Berbasis XGBoost Gain-based Importance}
Tahap ini bertujuan untuk mencapai efisiensi komputasi yang tinggi tanpa mengorbankan informasi krusial yang diperlukan untuk mendeteksi anomali. Mengingat dataset CSE-CIC-IDS2018 memiliki dimensi fitur yang besar (68 kolom setelah pra-pemrosesan), reduksi fitur berbasis \textit{Gain-based Importance} dilakukan melalui tahapan teknis sistematis berikut:

\begin{enumerate}
    \item \textbf{Pelatihan Model Awal (\textit{Baseline Training}):} Tahap awal diawali dengan melatih model XGBoost \textit{baseline} menggunakan seluruh fitur yang tersedia ($68$ fitur). Dataset dibagi menjadi subset data latih (\textit{training set}) dan data uji (\textit{test set}) menggunakan rasio $80:20$.
    
    \item \textbf{Ekstraksi Metrik \textit{Gain}:} Setelah proses pelatihan \textit{baseline} selesai, metrik \textit{Gain} diekstraksi dari fungsi internal XGBoost. Metrik \textit{Gain} merepresentasikan kontribusi rata-rata dari suatu fitur dalam menurunkan fungsi kerugian (\textit{loss function}) pada setiap pemisahan (\textit{split}) di percabangan pohon keputusan. Fitur dengan nilai \textit{Gain} tertinggi mengindikasikan kekuatan diskriminatif paling dominan dalam membedakan trafik normal (\textit{benign}) dari berbagai kategori serangan.
    
    \item \textbf{Perangkaian dan Seleksi Fitur (\textit{Feature Ranking}):} Seluruh fitur diurutkan dari nilai \textit{Gain} tertinggi hingga terendah. Fitur dengan skor \textit{Gain} paling tinggi menunjukkan kontribusi terbesar dalam pemisahan kelas pada pohon keputusan. Hasil ranking lengkap Top-20 fitur disajikan pada Bagian 3.2 (Tabel \ref{tab:feature_importance}).

    \item \textbf{Validasi melalui \textit{Ablation Study}:} Untuk menentukan batas jumlah fitur optimal, dilakukan \textit{Ablation Study} dengan melatih ulang model pada beberapa konfigurasi subset fitur, yaitu: \textit{Full Features} (68), Top-15, Top-10, dan Top-5. Performa setiap subset dievaluasi menggunakan metrik $F_1\text{-score}$ dan MCC, serta dibandingkan terhadap model \textit{baseline} berdimensi penuh.
    
    \item \textbf{Kriteria Penentuan Konfigurasi Optimal:} Pemilihan konfigurasi fitur terbaik didasarkan pada keseimbangan antara tiga aspek: (a) retensi performa klasifikasi ($F_1\text{-score}$ dan MCC), (b) efisiensi ukuran artefak model (\textit{model size}), dan (c) kecepatan inferensi (\textit{inference latency}). Konfigurasi yang memenuhi ketiga aspek secara berimbang ditetapkan sebagai titik keseimbangan (\textit{sweet spot}). Hasil lengkap validasi ini disajikan pada Bagian \ref{sec:hasil_ablation}.
\end{enumerate}

% 2.3
\subsection{Pembangkitan Sampel Adversarial (\textit{Evasion Attack})}
Tahap ini bertujuan untuk mensimulasikan skenario penyerang yang berusaha mengelabui model NIDS melalui serangan \textit{evasion}. Pembangkitan sampel \textit{adversarial} memanfaatkan teknik berbasis gradien matematis untuk mengidentifikasi kerentanan model dan memberikan perturbasi halus pada fitur-fitur kritis.

\subsubsection{Landasan Matematis \textit{Saliency Map}}
Misalkan model NIDS didefinisikan sebagai fungsi keputusan $f(\mathbf{x})$ dengan parameter $\theta$, di mana $\mathbf{x} = [x_1, x_2, \dots, x_d]^T \in \mathbb{R}^d$ merepresentasikan vektor fitur trafik jaringan dan $y \in \{0, 1\}$ adalah label kebenaran (\textit{ground truth}). Tingkat sensitivitas fungsi kerugian (\textit{loss function}) $L(\theta, \mathbf{x}, y)$ terhadap perubahan setiap elemen fitur diukur melalui vektor gradien:

\begin{equation}
\nabla_{\mathbf{x}} L(\theta, \mathbf{x}, y) = \left[ \frac{\partial L}{\partial x_1}, \frac{\partial L}{\partial x_2}, \dots, \frac{\partial L}{\partial x_d} \right]^T
\end{equation}

Skor \textit{Saliency Map} $S(\mathbf{x}, i)$ untuk fitur ke-$i$ dirumuskan sebagai nilai mutlak turunan parsial \textit{loss} terhadap fitur tersebut:

\begin{equation}
S(\mathbf{x}, i) = \left| \frac{\partial L(\theta, \mathbf{x}, y)}{\partial x_i} \right|
\end{equation}

Fitur yang memiliki nilai $S(\mathbf{x}, i)$ tinggi menunjukkan bahwa perubahan nilai secara minimal pada fitur tersebut akan memberikan pengaruh paling signifikan terhadap pergeseran output keputusan model.

\subsubsection{Formulasi Perturbasi Epsilon ($\epsilon$)}
Setelah fitur-fitur paling sensitif diidentifikasi, sampel \textit{adversarial} $\mathbf{x}_{adv}$ dibangkitkan dengan menambahkan vektor gangguan (\textit{perturbation}) $\boldsymbol{\delta}$ ke sampel asli $\mathbf{x}$. Untuk menjaga agar karakteristik trafik tetap fungsional dan tidak memicu anomali ekstrim, besaran gangguan dibatasi oleh batasan norma $L_\infty$ bernilai $\epsilon$:

\begin{equation}
\mathbf{x}_{adv} = \mathbf{x} + \boldsymbol{\delta}, \quad \text{dengan } \|\boldsymbol{\delta}\|_\infty \le \epsilon
\end{equation}

Secara eksplisit, proses pembentukan sampel \textit{adversarial} mengikuti formulasi arah gradien:

\begin{equation}
\mathbf{x}_{adv} = \mathbf{x} + \epsilon \cdot \text{sign}\left( \nabla_{\mathbf{x}} L(\theta, \mathbf{x}, y) \right)
\end{equation}

di mana $\text{sign}(\cdot)$ adalah fungsi tanda yang mengarahkan pergeseran fitur ke arah peningkatan kerugian model secara maksimal.

\subsubsection{Keterkaitan dengan Model XGBoost Tereduksi}
Pada model NIDS yang telah dioptimasi fiturnya melalui seleksi berbasis metrik \textit{Gain} (prosedur seleksi dijelaskan pada Bagian 2.2), konsentrasi keputusan model berada sepenuhnya pada dimensi yang sangat terbatas. Dalam penelitian ini, konfigurasi fitur tereduksi (yang ditentukan berdasarkan hasil \textit{Ablation Study} pada Bagian 3.3) digunakan sebagai objek pengujian ketahanan terhadap serangan \textit{evasion}. Secara matematis, semakin sedikit fitur yang digunakan, maka nilai gradien $\left| \frac{\partial L}{\partial x_i} \right|$ pada fitur-fitur tersebut menjadi jauh lebih besar dibandingkan pada model dengan fitur penuh ($68$ fitur), karena beban diskriminasi terkonsentrasi pada dimensi yang lebih sempit.

Kelemahan matematis ini dieksploitasi melalui simulasi serangan \textit{Saliency Map} berbasis $\epsilon$ untuk menghasilkan himpunan data $\mathbf{x}_{adv}$. Data terperturbasi ini selanjutnya dijadikan input pada fase \textit{Adversarial Training} guna memperkuat batas keputusan (\textit{decision boundary}) model terhadap manipulasi fitur halus di lingkungan produksi.


% 2.4
\subsection{Modelling dan Adversarial Training}
Fase ini berfokus pada perancangan arsitektur model klasifikasi XGBoost serta implementasi mekanisme pertahanan berbasis \textit{Adversarial Training}. Pendekatan ini bertujuan untuk merestrukturisasi batas keputusan (\textit{decision boundary}) model agar tidak mudah ditembus oleh sampel \textit{evasion} terperturbasi.

% --- ILUSTRASI DIAGRAM ALUR KERJA ADVERSARIAL TRAINING ---
\begin{figure}[htbp]
\centering
\vspace{0.2cm}
\resizebox{0.95\linewidth}{!}{%
\begin{tikzpicture}[
    node distance=1.2cm and 1.5cm,
    block/.style={rectangle, draw=blue!70, fill=blue!5, thick, rounded corners, align=center, minimum height=0.8cm, font=\small},
    process/.style={rectangle, draw=orange!80, fill=orange!10, thick, rounded corners, align=center, minimum height=0.8cm, font=\small},
    model/.style={ellipse, draw=green!60!black, fill=green!10, thick, align=center, minimum height=0.8cm, font=\small},
    arrow/.style={-Stealth, thick, draw=gray!80}
]
    % Nodes
    \node (clean) [block] {Dataset Latih Clean\\($\mathcal{D}_{clean}$)};
    \node (saliency) [process, below=of clean] {Saliency Map \& Perturbasi $\epsilon$\\$\mathbf{x}_{adv} = \mathbf{x} + \epsilon \cdot \text{sign}(\nabla_\mathbf{x} L)$};
    \node (adv_data) [block, below=of saliency] {Sampel Adversarial\\($\mathcal{D}_{adv}$)};
    \node (combine) [process, right=1.2cm of saliency] {Augmentasi Data\\$\mathcal{D}_{robust} = \mathcal{D}_{clean} \cup \mathcal{D}_{adv}$\\(Rasio $80:20$)};
    \node (retrain) [model, right=1.2cm of combine] {Pelatihan Ulang XGBoost\\(\textit{Min-Max Optimization})};
    \node (robust_model) [block, draw=green!70!black, fill=green!5, below=of retrain] {Robust NIDS Model\\(Model Terpengaruh)};

    % Arrows
    \draw [arrow] (clean) -- (saliency);
    \draw [arrow] (saliency) -- (adv_data);
    \draw [arrow] (clean.east) -- ++(0.6,0) |- (combine.north west);
    \draw [arrow] (adv_data.east) -- ++(0.6,0) |- (combine.south west);
    \draw [arrow] (combine) -- (retrain);
    \draw [arrow] (retrain) -- (robust_model);
\end{tikzpicture}%
}
\vspace{0.2cm}
% --- KETERANGAN GAMBAR PINDAH KE BAWAH GAMBAR ---
\caption{Diagram Alur Kerja Sistematis \textit{Adversarial Training}}
\label{fig:adversarial_training_flow}
\end{figure}



Alur kerja mekanisme \textit{Adversarial Training} secara menyeluruh diilustrasikan pada Gambar \ref{fig:adversarial_training_flow}, yang mencakup pembentukan sampel terperturbasi hingga pelatihan ulang model.

\subsubsection{Konfigurasi Model Klasifikasi XGBoost}
Model NIDS dibangun menggunakan algoritma \textit{Extreme Gradient Boosting} (XGBoost) yang dioptimasi pada subset fitur tereduksi hasil seleksi metrik \textit{Gain} (konfigurasi optimal ditentukan melalui \textit{Ablation Study} pada Bagian 3.3). Fungsi tujuan (\textit{objective function}) yang digunakan adalah \texttt{multi:softprob} untuk mengklasifikasikan trafik ke dalam kelas \textit{benign} dan berbagai variasi serangan jaringan. Hiperparameter utama disesuaikan untuk mencegah \textit{overfitting}, dengan kedalaman pohon maksimum (\texttt{max\_depth}) sebesar $6$, tingkat pembelajaran (\texttt{learning\_rate}) $0.1$, dan jumlah estimator (\texttt{n\_estimators}) $100$.

\subsubsection{Formulasi Matematis \textit{Adversarial Training}}
\textit{Adversarial Training} diformulasikan sebagai masalah optimasi \textit{Min-Max} (\textit{Saddle-Point Optimization}). Secara intuitif, tujuannya adalah meminimalkan kerugian maksimum (\textit{worst-case loss}) yang diakibatkan oleh perturbasi penyerang:

\begin{equation}
\min_{\theta} \mathbb{E}_{(\mathbf{x}, y) \sim \mathcal{D}} \left[ \max_{\|\boldsymbol{\delta}\|_\infty \le \epsilon} L(\theta, \mathbf{x} + \boldsymbol{\delta}, y) \right]
\end{equation}

di mana $\theta$ merepresentasikan parameter bobot model XGBoost, $\mathcal{D}$ adalah distribusi data asli, $L(\cdot)$ adalah fungsi kerugian \textit{cross-entropy}, dan $\boldsymbol{\delta}$ adalah vektor perturbasi \textit{Saliency Map} yang dibatasi oleh norma $\|\boldsymbol{\delta}\|_\infty \le \epsilon$.

\subsubsection{Strategi Augmentasi dan Pelatihan Ulang (\textit{Re-training})}
Implementasi praktis dari fungsi tujuan \textit{Min-Max} dilakukan melalui teknik augmentasi dataset latih dengan alur sebagai berikut:
\begin{enumerate}
    \item \textbf{Generasi Sampel Pertahanan:} Menggunakan dataset latih bersih ($\mathcal{D}_{clean}$), dibangkitkan subset sampel \textit{adversarial} ($\mathcal{D}_{adv}$) menggunakan metode \textit{Saliency Map} dengan variasi $\epsilon \in \{0.01, 0.1\}$.
    \item \textbf{Konstruksi Dataset Gabungan:} Dataset latih baru ($\mathcal{D}_{robust}$) dibentuk dengan menggabungkan data murni dan sampel \textit{adversarial} dengan proporsi rasio $80:20$:
    \begin{equation}
    \mathcal{D}_{robust} = \mathcal{D}_{clean} \cup \mathcal{D}_{adv}
    \end{equation}
    \item \textbf{Restrukturisasi Batas Keputusan:} Model XGBoost dilatih ulang pada $\mathcal{D}_{robust}$. Penambahan sampel terperturbasi memicu penyesuaian ambang batas pemisahan (\textit{split threshold}) pada cabang-cabang pohon keputusan. Akibatnya, wilayah sekitar sampel asli yang semula rentan terhadap pergeseran $\epsilon$ kini tercakup di dalam area keputusan yang benar, sehingga meningkatkan ketangguhan (\textit{robustness}) model secara signifikan tanpa mengorbankan akurasi pada trafik normal.
\end{enumerate}



% 2.5
\subsection{Robustness Ablation Study}
\textit{Robustness Ablation Study} dirancang sebagai tahap analisis komparatif mendalam untuk mengevaluasi trade-off matematis antara efisiensi dimensi fitur dan tingkat keamanan model NIDS terhadap serangan \textit{evasion}.

\subsubsection{Skenario Konfigurasi Subset Fitur}
Pengujian dilakukan dengan mengisolasi variabel fitur ke dalam empat tingkatan hirarki berbasis peringkat metrik \textit{Gain} XGBoost, sebagaimana dirinci pada Tabel \ref{tab:ablation_configurations}.

% --- Pastikan \usepackage{makecell} dan \usepackage{booktabs} sudah ada di Preamble ---

\begin{table}[htbp]
\caption{Konfigurasi Fitur untuk Robustness Ablation Study}
\label{tab:ablation_configurations}
\centering
\vspace{0.2cm}
\small
\begin{tabular}{@{}c c p{0.55\linewidth}@{}}
\toprule
\textbf{Konfigurasi} & \textbf{Jumlah} & \textbf{Deskripsi} \\ \midrule
$C_1$ (\textit{Full}) & 68 & Basis komparasi dimensi penuh \\ \addlinespace
$C_2$ (\textit{Top-15}) & 15 & Subset fitur informatif moderat \\ \addlinespace
$C_3$ (\textit{Top-10}) & 10 & Hipotesis \textit{sweet spot} (ringkas \& presisi) \\ \addlinespace
$C_4$ (\textit{Top-5}) & 5 & Kompresi ekstrim (efisiensi maksimum) \\ \bottomrule
\end{tabular}
\end{table}


\subsubsection{Kriteria Evaluasi dan Titik Keseimbangan (\textit{Sweet Spot})}
Setiap konfigurasi $C_k$ ($k \in \{1, 2, 3, 4\}$) dilatih menggunakan dua pendekatan, yakni model \textit{baseline} (tanpa pertahanan) dan model \textit{robust} (dengan \textit{Adversarial Training}), kemudian diuji menggunakan sampel \textit{evasion} $\mathbf{x}_{adv}$ berparameter $\epsilon$. Tujuannya adalah mengidentifikasi titik pembanding matematis (\textit{sweet spot}) $\hat{C}$ yang memaksimalkan fungsi utilitas gabungan:

\begin{equation}
\hat{C} = \arg\max_{C_k} \left( \text{MCC}(C_k, \mathbf{x}_{adv}) - \lambda \cdot |C_k| \right)
\end{equation}

di mana $\text{MCC}(C_k, \mathbf{x}_{adv})$ mengukur ketangguhan deteksi pada trafik yang dimanipulasi, $|C_k|$ melambangkan jumlah fitur, dan $\lambda$ adalah bobot penalti efisiensi dimensi. Hipotesis awal penelitian ini menetapkan bahwa terdapat titik keseimbangan pada salah satu konfigurasi $C_k$ yang mampu memangkas ukuran model secara signifikan tanpa mengorbankan ketahanan terhadap serangan \textit{evasion}. Pembuktian empiris hipotesis ini disajikan pada Bagian 3.6.



% 2.6
\subsection{Evaluasi Model dan Perbandingan Performa}
Untuk menilai efektivitas model secara menyeluruh, evaluasi dilakukan melalui skenario pengujian silang $2\times 2$ yang membandingkan model standar (\textit{baseline}) dengan model yang diperkuat (\textit{robust}) pada kondisi trafik murni maupun termanipulasi.

\subsubsection{Matriks Skenario Pengujian $2\times 2$}
Skenario pengujian dirancang untuk mengisolasi kerentanan model serta membuktikan efektivitas \textit{Adversarial Training}, sebagaimana disajikan pada Tabel \ref{tab:skenario_pengujian_2x2}.

% --- TABEL 1 KOLOM PENUH (AMAN & BEBAS ERROR) ---
\begin{table*}[t]
\caption{Matriks Skenario Evaluasi Model NIDS ($2\times 2$)}
\label{tab:skenario_pengujian_2x2}
\centering
\vspace{0.2cm}
\small
\begin{tabular}{@{} p{0.10\textwidth} p{0.15\textwidth} p{0.14\textwidth} p{0.14\textwidth} p{0.37\textwidth} @{}}
\toprule
\textbf{Skenario} & \textbf{Model} & \textbf{Data Latih} & \textbf{Data Uji} & \textbf{Tujuan Evaluasi} \\ \midrule
S1 & \textit{Baseline} & Murni & Murni & Performa awal model pada kondisi normal \\ \addlinespace
S2 & \textit{Vulnerability} & Murni & Manipulasi & Mengukur tingkat kerentanan \textit{evasion} (\textit{security gap}) \\ \addlinespace
S3 & \textit{Integritas} & \textit{Robust} & Murni & Memastikan akurasi model tetap terjaga pada trafik normal \\ \addlinespace
S4 & \textit{Robustness} & \textit{Robust} & Manipulasi & Membuktikan peningkatan ketahanan terhadap serangan \\ \bottomrule
\end{tabular}
\end{table*}

\subsubsection{Metrik Evaluasi Utama: Matthews Correlation Coefficient (MCC)}
Dataset NIDS seperti CSE-CIC-IDS2018 memiliki karakteristik ketidakseimbangan kelas (\textit{class imbalance}) yang sangat ekstrem antara trafik normal (\textit{benign}) dan serangan. Oleh karena itu, \textit{Matthews Correlation Coefficient} (MCC) dipilih sebagai metrik evaluasi utama. Berbeda dengan akurasi konvensional yang dapat memberikan gambaran bias pada data tidak seimbang, MCC memperhitungkan keempat kuadran \textit{Confusion Matrix} secara proporsional. Formulasi matematis MCC didefinisikan sebagai:

\begin{equation}
\text{MCC} = \frac{(TP \times TN) - (FP \times FN)}{\sqrt{(TP + FP)(TP + FN)(TN + FP)(TN + FN)}}
\end{equation}

di mana $TP$ adalah \textit{True Positive}, $TN$ adalah \textit{True Negative}, $FP$ adalah \textit{False Positive}, dan $FN$ adalah \textit{False Negative}. Nilai MCC berada pada rentang $[-1, +1]$, di mana $+1$ merepresentasikan prediksi sempurna, $0$ menunjukkan prediksi acak (\textit{random guessing}), dan $-1$ mengindikasikan ketidaksesuaian total antara prediksi dan kenyataan.

\subsubsection{Metrik Evaluasi Pendukung}
Sebagai pelengkap metrik MCC, performa model juga diukur menggunakan metrik $Precision$ dan $F_1\text{-score}$ untuk mengevaluasi ketepatan serta keseimbangan deteksi:

\begin{equation}
\text{Precision} = \frac{TP}{TP + FP}
\end{equation}

\begin{align}
F_1\text{-score} &= 2 \times \frac{\text{Precision} \times \text{Recall}}{\text{Precision} + \text{Recall}} \\[6pt]
\text{dengan } \text{Recall} &= \frac{TP}{TP + FN}
\end{align}

Melalui kombinasi metrik ini, penurunan performa pada S2 menunjukkan celah keamanan (\textit{security gap}), sedangkan pemulihan nilai MCC dan $F_1\text{-score}$ pada S4 membuktikan keberhasilan penguatan model melalui \textit{Adversarial Training}.



% Bab 3
\section{HASIL DAN PEMBAHASAN}


% 3.1
\subsection{Hasil Pra-pemrosesan Data dan Deskripsi Subset}
Tahap pra-pemrosesan data bertujuan untuk mentransformasikan dataset mentah CSE-CIC-IDS2018 yang berukuran masif ($\approx 6{,}7\text{ GB}$) menjadi subset data yang bersih, seimbang secara proporsional, dan siap digunakan untuk pelatihan model.

\subsubsection{Reduksi Volume dan \textit{Stratified Sampling}}
Mengingat keterbatasan memori komputasi serta untuk mempercepat iterasi eksperimen, teknik \textit{Stratified Sampling} sebesar 10\% diterapkan pada seluruh file CSV dataset. Penerapan stratifikasi ini menjamin bahwa rasio kemunculan setiap kategori serangan asli tetap dipertahankan secara teratur. Tabel \ref{tab:distribusi_dataset} menyajikan distribusi sampel data sebelum dan sesudah proses pembersihan serta penerapan sampling.

% --- TABEL 1 KOLOM (SINGLE COLUMN FIT) ---
\begin{table}[htbp]
\caption{Distribusi Sampel Data CSE-CIC-IDS2018 Sebelum dan Sesudah Pra-pemrosesan}
\label{tab:distribusi_dataset}
\centering
\vspace{0.2cm}
\scriptsize
\begin{tabular}{@{} l r r p{0.30\linewidth} @{}}
\toprule
\textbf{Kategori} & \textbf{Data Awal} & \textbf{Subset 10\%} & \textbf{Deskripsi} \\ \midrule
\textit{Benign} & 13.484.708 & 1.348.470 & Trafik normal \\ \addlinespace
\textit{DDOS / DOS} & 1.391.082 & 139.108 & \textit{Flooding} volumetrik \\ \addlinespace
\textit{Brute-Force} & 380.943 & 38.094 & \textit{Login} SSH/FTP \\ \addlinespace
\textit{Botnet} & 286.191 & 28.619 & Komando \textit{C2} \\ \addlinespace
\textit{Web Attacks} & 928 & 93 & \textit{SQLi / XSS} \\ \addlinespace
\textit{Infiltration} & 161.934 & 16.193 & Penetrasi internal \\ \midrule
\textbf{Total} & \textbf{15.705.786} & \textbf{1.570.577} & 78 Fitur (sebelum cleaning) \\ \bottomrule
\end{tabular}
\end{table}

\subsubsection{Hasil Pembersihan Data (\textit{Data Cleaning})}

\begin{figure}[htbp]
\centering
\includegraphics[width=0.85\linewidth]{data/bar_distribusi_label_populasi.png}
\caption{Distribusi Label Populasi Asli pada Dataset CSE-CIC-IDS2018}
\label{fig:distribusi_label}
\end{figure}

Proses \textit{data cleaning} secara sistematis memberikan dampak langsung pada kualitas dan integritas dataset:
\begin{enumerate}
    \item \textbf{Penanganan Nilai Singular:} Sebanyak $2.861$ baris sampel yang mengandung nilai tak terdefinisi (\textit{NaN/Missing}) berhasil dieliminasi. Nilai \textit{Infinity} ($\infty$) pada variabel kuantitatif seperti \texttt{Flow Bytes/s} diimputasi menggunakan nilai terbatas maksimum (\textit{finite max}) pada kolom masing-masing untuk menjaga kestabilan matematis saat kalkulasi gradien.
    \item \textbf{Pembersihan Fitur Bias dan Variansi Nol:} Penghapusan $6$ fitur identitas aliran (\textit{Flow ID}, \textit{Source/Destination IP}, \textit{Ports}, \textit{Protocol}) serta variabel \textit{Timestamp} berhasil mencegah model mengalami \textit{overfitting} terhadap konfigurasi topologi jaringan spesifik. Pemindaian variansi menghasilkan pembuangan $8$ fitur yang bersifat konstan (\textit{zero-variance}).
    \item \textbf{Matriks Akhir Fitur:} Dari total $83$ fitur teknis awal, proses pra-pemrosesan menghasilkan matriks data bersih (\texttt{cleaned\_100.pkl}) dengan $68$ fitur numerik independen dan $1$ kolom label target, yang siap dimasukkan ke dalam tahap perangkaian fitur XGBoost \textit{Gain}.
\end{enumerate}


% 3.2
\subsection{Analisis Ranking Fitur XGBoost Gain dan Eksplorasi Fitur Sensitif}
Tahap ini menguraikan hasil ekstraksi kepentingan fitur (\textit{feature importance}) berbasis metrik \textit{Gain} pada model XGBoost \textit{baseline}, serta menganalisis perbedaan mendasar antara fitur penting secara global (\textit{global importance}) dan fitur sensitif secara lokal (\textit{local sensitivity}) melalui pendekatan \textit{Saliency Map}.

\subsubsection{Distribusi Kontribusi Metrik \textit{Gain} XGBoost}
Hasil pelatihan model \textit{baseline} menggunakan seluruh $68$ fitur menunjukkan bahwa distribusi kontribusi \textit{Gain} tidak tersebar merata, melainkan terkonsentrasi secara dominan pada beberapa atribut utama. Tabel \ref{tab:feature_importance} menyajikan daftar 20 fitur teratas berdasarkan skor \textit{Gain} yang diekstraksi dari model XGBoost \textit{baseline}.

% --- TABEL TOP 20 FEATURES (DIPINDAH DARI BAB 2) ---
\begin{table}[htbp]
\caption{Top 20 Features Ranked by XGBoost Gain-based Importance}
\label{tab:feature_importance}
\centering
\vspace{0.2cm}
\small
\begin{tabular}{@{}cl@{}}
\toprule
\textbf{Rank} & \textbf{Feature Name} \\ \midrule
1 & Fwd Seg Size Min \\
2 & URG Flag Cnt \\
3 & Tot Bwd Pkts \\
4 & Fwd Act Data Pkts \\
5 & Fwd Pkt Len Max \\
6 & Fwd Pkt Len Mean \\
7 & Bwd Pkt Len Mean \\
8 & Init Bwd Win Byts \\
9 & TotLen Bwd Pkts \\
10 & Init Fwd Win Byts \\ \midrule
11 & Bwd Pkt Len Std \\
12 & Bwd Seg Size Avg \\
13 & Fwd IAT Std \\
14 & Bwd IAT Max \\
15 & Bwd Pkt Len Max \\
16 & Fwd Header Len \\
17 & Fwd Pkts/s \\
18 & Fwd URG Flags \\
19 & Fwd IAT Tot \\
20 & Fwd Pkt Len Std \\ \bottomrule
\end{tabular}
\end{table}

Fitur-fitur seperti \textit{Fwd Seg Size Min}, \textit{Init Bwd Win Byts}, dan \textit{Tot Bwd Pkts} menduduki peringkat teratas karena merepresentasikan variabel kontrol protokol TCP/IP yang sangat spesifik dalam membedakan pola inisiasi koneksi.

Secara kumulatif, sepuluh fitur teratas (Top-10) menyumbang lebih dari $85\%$ dari total akumulasi skor \textit{Gain} seluruh model. Temuan ini membuktikan secara empiris bahwa mayoritas dari $68$ fitur bersifat redundan atau memberikan informasi diskriminatif yang sangat minim bagi algoritma pemutus keputusan. Visualisasi distribusi kontribusi \textit{Gain} disajikan pada Gambar \ref{fig:feature_importance_top20}.

\begin{figure}[htbp]
\centering
\includegraphics[width=0.85\linewidth]{data/feature_importance_top20.png}
\caption{Visualisasi Feature Importance Top-20 Berdasarkan XGBoost Gain Metric}
\label{fig:feature_importance_top20}
\end{figure}

\subsubsection{Komparasi \textit{Global Importance} vs. \textit{Local Sensitivity}}
Meskipun pemangkasan fitur ke Top-10 terbukti mempertahankan akurasi klasifikasi pada kondisi normal, analisis berbasis gradien mengungkapkan dinamika kerentanan baru pada tingkat sampel individual:
\begin{enumerate}
    \item \textbf{\textit{Global Importance} (XGBoost \textit{Gain}):} Mengukur kontribusi rata-rata suatu fitur dalam mereduksi fungsi kerugian (\textit{loss function}) di seluruh percabangan pohon keputusan selama fase pelatihan. Fitur dengan nilai \textit{Gain} tinggi menjadi tumpuan utama model dalam melakukan inferensi umum pada data bersih.
    
    \item \textbf{\textit{Local Sensitivity} (\textit{Saliency Map Gradient}):} Mengukur besarnya nilai mutlak turunan parsial $\left| \frac{\partial L}{\partial x_i} \right|$ untuk sampel tertentu $x_i$. Evaluasi gradien menunjukkan bahwa ketika model dipaksa mengambil keputusan hanya dari $10$ fitur, besaran gradien pada $10$ fitur tersebut meningkat secara drastis dibandingkan pada model $68$ fitur.
\end{enumerate}

\subsubsection{Implikasi Terhadap Celah Keamanan (\textit{Evasion Vulnerability})}
Peningkatan kecuraman gradien $\left| \frac{\partial L}{\partial x_i} \right|$ pada subset Top-10 fitur (khususnya pada atribut \textit{Init Bwd Win Byts} dan \textit{Flow Pkts/s}) menunjukkan bahwa model terreduksi menjadi sangat reaktif terhadap pergeseran nilai pada variabel-variabel tersebut. Kondisi matematis ini memberikan ruang bagi penyerang untuk mengeksploitasi model: dengan mengaplikasikan perturbasi $\epsilon$ bernilai sangat kecil ($\epsilon = 0.01$) pada fitur terisolasi tersebut, penyerang mampu memicu lonjakan \textit{loss} yang cukup besar untuk membelokkan prediksi model dari klasifikasi serangan menjadi trafik normal (\textit{benign}).

\begin{figure}[htbp]
\centering
\includegraphics[width=0.85\linewidth]{data/saliency_map_top10.png}
\caption{Saliency Map Top-10 Fitur: Sensitivitas Model terhadap Perubahan Fitur}
\label{fig:saliency_map}
\end{figure}

Gambar \ref{fig:saliency_map} memvisualisasikan skor \textit{saliency} (nilai gradien absolut rata-rata $\left|\frac{\partial L}{\partial x_i}\right|$) untuk masing-masing dari 10 fitur terpilih. Semakin tinggi nilai saliency suatu fitur, semakin besar dampak perubahan kecil pada fitur tersebut terhadap keputusan model. Temuan utama dari visualisasi ini adalah bahwa fitur \textit{Init Fwd Win Byts} (Initial Forward Window Bytes) memiliki skor saliency tertinggi, menjadikannya fitur paling sensitif terhadap perturbasi. Menariknya, fitur ini berada pada peringkat ke-10 dalam metrik \textit{Gain} (Tabel \ref{tab:feature_importance}) --- bukan fitur yang paling penting saat pelatihan, namun justru yang paling rentan saat diserang. Secara teknis, \textit{Init Fwd Win Byts} merepresentasikan ukuran \textit{TCP window size} awal dari koneksi \textit{forward}. Penyerang dapat memanipulasi nilai ini pada header TCP tanpa merusak fungsionalitas paket jaringan, sehingga menjadikannya target serangan \textit{evasion} yang sangat realistis. Perbedaan antara peringkat \textit{Gain} dan skor \textit{saliency} ini mengkonfirmasi bahwa fitur yang paling penting secara global saat pelatihan tidak selalu identik dengan fitur yang paling rentan secara lokal saat inferensi.

\begin{figure}[htbp]
\centering
\includegraphics[width=0.85\linewidth]{data/vulnerability_mcc_f1_vs_epsilon.png}
\caption{Kurva Degradasi MCC dan F1-Score terhadap Variasi Epsilon Perturbasi}
\label{fig:vulnerability_curve}
\end{figure}

Gambar \ref{fig:vulnerability_curve} mengilustrasikan dampak peningkatan intensitas perturbasi $\epsilon$ terhadap penurunan metrik MCC dan $F_1$-score pada model \textit{baseline} (tanpa pertahanan). Terlihat bahwa pada $\epsilon = 0{,}01$, model sudah mulai mengalami degradasi yang terukur. Seiring kenaikan $\epsilon$ ke $0{,}05$ dan $0{,}10$, kedua metrik turun secara progresif hingga mendekati nol untuk MCC. Perbedaan perilaku antara MCC dan $F_1$-score juga terlihat jelas: $F_1$-score tidak turun sedrastis MCC karena model masih berhasil mengklasifikasikan kelas mayoritas (\textit{benign}) dengan benar, namun gagal total pada kelas serangan. MCC yang memperhitungkan seluruh kuadran \textit{Confusion Matrix} secara berimbang mampu menangkap kolaps ini dengan lebih sensitif. Kurva degradasi ini membuktikan urgensi penerapan mekanisme pertahanan (\textit{Adversarial Training}) yang dibahas pada Bagian 3.5.


% 3.3
\subsection{Hasil Ablation Study Seleksi Fitur}
\label{sec:hasil_ablation}
Berdasarkan prosedur \textit{Ablation Study} yang dijelaskan pada Bagian 2.2, model XGBoost dilatih ulang pada empat konfigurasi subset fitur. Tabel \ref{tab:ablation_feature_selection} menyajikan perbandingan performa klasifikasi, ukuran artefak model, dan kecepatan inferensi untuk setiap konfigurasi.

% --- TABEL HASIL ABLATION STUDY SELEKSI FITUR ---
\begin{table}[htbp]
\caption{Hasil Ablation Study Seleksi Fitur: Perbandingan Performa dan Efisiensi}
\label{tab:ablation_feature_selection}
\centering
\vspace{0.2cm}
\resizebox{\linewidth}{!}{%
\begin{tabular}{@{} l c c c c c @{}}
\toprule
\textbf{Konfigurasi} & \textbf{Fitur} & \textbf{$F_1$-score} & \textbf{MCC} & \textbf{Ukuran Model} & \textbf{Inferensi (10k)} \\ \midrule
\textit{Full Features} & 68 & 0.9780 & 0.9510 & 612 KB & 0.2400 s \\ \addlinespace
Top-15 & 15 & 0.9752 & 0.9442 & 124 KB & 0.1950 s \\ \addlinespace
\textbf{Top-10} & \textbf{10} & \textbf{0.9728} & \textbf{0.9351} & \textbf{84 KB} & \textbf{0.1826 s} \\ \addlinespace
Top-5 & 5 & 0.9210 & 0.8810 & 42 KB & 0.1520 s \\ \bottomrule
\end{tabular}%
}
\end{table}

Berdasarkan Tabel \ref{tab:ablation_feature_selection}, beberapa temuan kuantitatif dapat dianalisis sebagai berikut:
\begin{enumerate}
    \item \textbf{Konfigurasi Full Features (68 fitur)} mencatatkan performa tertinggi dengan $F_1 = 0{,}9780$ dan $\text{MCC} = 0{,}9510$. Namun, ukuran model mencapai $612\text{ KB}$ dan waktu inferensi $0{,}2400\text{ s}$ per 10.000 sampel --- yang kurang efisien untuk skenario \textit{real-time}.
    \item \textbf{Konfigurasi Top-15 (15 fitur)} mengalami penurunan performa yang sangat marginal ($\Delta F_1 = 0{,}29\%$; $\Delta\text{MCC} = 0{,}72\%$) dengan pengurangan ukuran model sebesar $80\%$ (dari $612\text{ KB}$ ke $124\text{ KB}$).
    \item \textbf{Konfigurasi Top-10 (10 fitur)} menunjukkan penurunan $F_1$ hanya $0{,}53\%$ dari \textit{Full Features}, sementara ukuran model berkurang $86\%$ menjadi $84\text{ KB}$ dan waktu inferensi turun $24\%$ menjadi $0{,}1826\text{ s}$.
    \item \textbf{Konfigurasi Top-5 (5 fitur)} mengalami degradasi performa yang signifikan ($\Delta F_1 = 5{,}83\%$; $\Delta\text{MCC} = 7{,}36\%$), menandakan bahwa 5 fitur tidak memiliki kapasitas diskriminatif yang cukup untuk membedakan seluruh kategori serangan.
\end{enumerate}

\subsubsection{Visualisasi Trade-off Performa vs. Efisiensi Seleksi Fitur}
Gambar \ref{fig:ablation_feature_tradeoff} mengilustrasikan kurva trade-off antara jumlah fitur, skor $F_1$, dan ukuran model. Terlihat bahwa penurunan performa dari Full (68) ke Top-10 sangat kecil ($\Delta F_1 < 0{,}6\%$), sedangkan efisiensi ukuran model meningkat drastis (reduksi $86\%$).

% --- GRAFIK ABLATION STUDY FEATURE SELECTION ---
\begin{figure}[htbp]
\centering
\vspace{0.2cm}
\resizebox{0.65\linewidth}{!}{%
\begin{tikzpicture}
\begin{axis}[
    xlabel={Jumlah Fitur},
    ylabel={$F_1$-score},
    axis y line*=left,
    xmin=3, xmax=72,
    ymin=0.90, ymax=1.00,
    xtick={5, 10, 15, 68},
    ytick={0.90, 0.92, 0.94, 0.96, 0.98, 1.00},
    grid=major,
    grid style={dashed, gray!30},
    width=9cm,
    height=6.5cm,
    thick,
    ylabel style={color=blue!70!black}
]
% Kurva F1-score (Sumbu Kiri)
\addplot[
    color=blue!70!black,
    mark=circle*,
    line width=1.5pt
] coordinates {
    (5, 0.9210)
    (10, 0.9728)
    (15, 0.9752)
    (68, 0.9780)
};
\label{plot_f1_ablation}
\end{axis}

\begin{axis}[
    ylabel={Ukuran Model (KB)},
    axis y line*=right,
    axis x line=none,
    xmin=3, xmax=72,
    ymin=0, ymax=700,
    xtick={5, 10, 15, 68},
    width=9cm,
    height=6.5cm,
    thick,
    ylabel style={color=red!70!black}
]
% Kurva Ukuran Model (Sumbu Kanan)
\addplot[
    color=red!70!black,
    mark=square*,
    line width=1.5pt,
    dashed
] coordinates {
    (5, 42)
    (10, 84)
    (15, 124)
    (68, 612)
};
\label{plot_size_ablation}
\end{axis}
\end{tikzpicture}%
}
\vspace{0.2cm}
\caption{Trade-off Seleksi Fitur: $F_1$-score (garis biru, sumbu kiri) vs. Ukuran Model (garis merah putus-putus, sumbu kanan). Konfigurasi Top-10 memberikan keseimbangan optimal.}
\label{fig:ablation_feature_tradeoff}
\end{figure}

\subsubsection{Penetapan Top-10 sebagai Konfigurasi Optimal}
Berdasarkan Tabel \ref{tab:ablation_feature_selection} dan Gambar \ref{fig:ablation_feature_tradeoff}, konfigurasi Top-10 fitur dipilih sebagai titik keseimbangan (\textit{sweet spot}) dengan pertimbangan:
\begin{enumerate}
    \item \textbf{Retensi Performa:} Penurunan $F_1$-score dari konfigurasi penuh (68 fitur) ke Top-10 hanya sebesar $0{,}53\%$ (dari $0{,}9780$ ke $0{,}9728$), yang secara praktis tidak signifikan.
    \item \textbf{Efisiensi Ukuran:} Ukuran model tereduksi $86\%$ (dari $612\text{ KB}$ menjadi $84\text{ KB}$), menjadikannya ideal untuk \textit{deployment} pada perangkat \textit{edge} dan arsitektur \textit{serverless}.
    \item \textbf{Kecepatan Inferensi:} Waktu inferensi per 10.000 sampel turun $24\%$ (dari $0{,}2400\text{ s}$ menjadi $0{,}1826\text{ s}$).
    \item \textbf{Batas Degradasi:} Konfigurasi Top-5 menunjukkan penurunan tajam ($\Delta F_1 = 5{,}7\%$; $\Delta\text{MCC} = 5{,}2\%$), mengkonfirmasi bahwa 5 fitur tidak memadai untuk mempertahankan performa klasifikasi.
\end{enumerate}

Konfigurasi Top-10 ini selanjutnya digunakan sebagai basis model untuk seluruh eksperimen \textit{Adversarial Training} dan \textit{Robustness Ablation} pada Bagian 3.4--3.7.


% 3.4
\subsection{Evaluasi Komparatif 4 Skenario Evaluasi Model (S1--S4)}
Untuk membuktikan efektivitas \textit{Adversarial Training} secara komprehensif, evaluasi dilakukan melintasi empat skenario pengujian silang $2\times 2$. Skenario ini membandingkan interaksi antara tipe model (\textit{Baseline} vs. \textit{Robust}) dengan tipe data uji (\textit{Murni} vs. \textit{Manipulasi}).

\subsubsection{Hasil Kuantitatif 4 Skenario Pengujian}
Tabel \ref{tab:results_4_scenarios} menyajikan hasil pengukuran kuantitatif lengkap untuk keempat skenario evaluasi pada subset fitur Top-10 dengan tingkat perturbasi $\epsilon = 0.10$.

% --- TABEL KUANTITATIF 4 SKENARIO (FIT 1 KOLOM) ---
% --- TABEL KUANTITATIF 4 SKENARIO (PASTI FIT 1 KOLOM) ---
\begin{table}[htbp]
\caption{Hasil Kuantitatif Evaluasi Model pada Skenario S1, S2, S3, dan S4}
\label{tab:results_4_scenarios}
\centering
\vspace{0.2cm}
\resizebox{\linewidth}{!}{%
\begin{tabular}{@{} p{0.32\linewidth} p{0.18\linewidth} p{0.18\linewidth} c c c @{}}
\toprule
\textbf{Skenario} & \textbf{Model} & \textbf{Data Uji} & \textbf{MCC} & \textbf{$F_1$-score} & \textbf{Precision} \\ \midrule
S1 (\textit{Baseline}) & \textit{Baseline} & Murni & 0.9351 & 0.9729 & 0.9720 \\ \addlinespace
S2 (\textit{Vulnerability}) & \textit{Baseline} & Manipulasi & 0.0184 & 0.7539 & 0.6905 \\ \addlinespace
S3 (\textit{Integritas}) & \textit{Robust} & Murni & 0.9347 & 0.9727 & 0.9718 \\ \addlinespace
S4 (\textit{Robustness}) & \textit{Robust} & Manipulasi & \textbf{0.9953} & \textbf{0.9985} & 0.9985 \\ \bottomrule
\end{tabular}%
}
\end{table}


\subsubsection{Visualisasi Perbandingan Performa 4 Skenario}
Perbandingan visual antara performa model standar dan model terfortifikasi pada seluruh skenario diilustrasikan pada Gambar \ref{fig:4_scenarios_comparison}.

% --- GRAFIK TIKZ 4 SKENARIO (FIT 1 KOLOM) ---
\begin{figure}[htbp]
\centering
\vspace{0.2cm}
\resizebox{0.60\linewidth}{!}{%
\begin{tikzpicture}
\begin{axis}[
    ybar,
    enlargelimits=0.15,
    legend style={at={(0.5,-0.25)}, anchor=north, legend columns=-1, font=\footnotesize},
    ylabel={Skor Metrik},
    ylabel style={font=\footnotesize},
    symbolic x coords={S1, S2, S3, S4},
    xtick=data,
    xticklabel style={font=\footnotesize},
    yticklabel style={font=\footnotesize},
    nodes near coords,
    nodes near coords align={vertical},
    every node near coord/.append style={font=\scriptsize},
    ymin=0, ymax=1.15,
    bar width=12pt,
    grid=major,
    grid style={dashed, gray!30}
]
\addplot[draw=blue!80, fill=blue!30] coordinates {(S1, 0.9351) (S2, 0.0184) (S3, 0.9347) (S4, 0.9953)};
\addplot[draw=orange!80, fill=orange!30] coordinates {(S1, 0.9729) (S2, 0.7539) (S3, 0.9727) (S4, 0.9985)};
\legend{MCC, $F_1$-score}
\end{axis}
\end{tikzpicture}%
}
\vspace{0.2cm}
\caption{Perbandingan Metrik MCC dan $F_1$-score pada Skenario S1, S2, S3, dan S4 ($\epsilon = 0.10$)}
\label{fig:4_scenarios_comparison}
\end{figure}

Gambar \ref{fig:4_scenarios_comparison} memperlihatkan kontras yang sangat tajam antara keempat skenario. Pada metrik MCC (bar biru), S2 nyaris menyentuh nol ($0{,}0184$) sementara ketiga skenario lainnya berada di atas $0{,}93$. Pola ini menunjukkan bahwa model \textit{baseline} mengalami kolaps total ketika berhadapan dengan data adversarial — namun masalah ini teratasi sepenuhnya setelah \textit{Adversarial Training} diterapkan (S4). Sementara itu, pada metrik $F_1$-score (bar oranye), S2 masih menunjukkan nilai $0{,}7539$ meskipun MCC-nya hampir nol. Fenomena ini terjadi karena $F_1$-score yang dihitung secara \textit{weighted average} masih terdongkrak oleh klasifikasi benar pada kelas mayoritas (\textit{benign}), sedangkan MCC yang memperhitungkan seluruh kuadran \textit{Confusion Matrix} lebih sensitif dalam mendeteksi kegagalan pada kelas minoritas (serangan).

\subsubsection{Analisis Dinamika Performa dan Keamanan}
Berdasarkan hasil eksperimen pada Tabel \ref{tab:results_4_scenarios} dan Gambar \ref{fig:4_scenarios_comparison}, temuan-temuan penting dapat dianalisis sebagai berikut:
\begin{enumerate}
    \item \textbf{Celah Keamanan S2 (\textit{Security Gap}):} Pada skenario S1, model \textit{baseline} berkinerja sangat tinggi ($\text{MCC} = 0{,}9351$). Namun, saat diuji dengan data terperturbasi pada S2, performanya anjlok hingga $\text{MCC} = 0{,}0184$ (penurunan 98\%). Hal ini membuktikan bahwa model standar sangat rentan terhadap manipulasi serangan \textit{evasion}.
    
    \item \textbf{Integritas Trafik Normal pada S3 (\textit{No Over-fitting/Paranoia}):} Pengujian model \textit{robust} pada data murni (S3) menghasilkan $\text{MCC} = 0{,}9347$, yang hampir identik dengan performa S1 (penurunan hanya 0,04\%). Temuan ini membuktikan bahwa penambahan sampel \textit{adversarial} tidak merusak kemampuan model dalam mengenali trafik jaringan normal.

    \item \textbf{Pemulihan Ketahanan pada S4 (\textit{Robustness Defense}):} Pada skenario S4, model yang dilatih ulang menggunakan \textit{Adversarial Training} mampu menahan paket yang dimanipulasi dengan mencapai nilai $\text{MCC} = 0{,}9953$. Peningkatan skor MCC sebesar $+0{,}9769$ dibanding S2 membuktikan bahwa batas keputusan (\textit{decision boundary}) baru berhasil mengisolasi pergeseran gradien serangan. Persentase pemulihan mencapai 106,6\% dari \textit{security gap}, yang berarti model \textit{robust} pada data adversarial justru melampaui performa model \textit{baseline} pada data murni.
    
    \item \textbf{Fenomena S4 $>$ S1 (\textit{Transfer Attack Effect}):} Nilai MCC S4 ($0{,}9953$) yang lebih tinggi dari S1 ($0{,}9351$) dapat dijelaskan oleh fakta bahwa sampel adversarial dibangkitkan menggunakan gradien model \textit{baseline} — bukan model \textit{robust}. Perturbasi yang dirancang untuk mengecoh \textit{baseline} menjadi tidak efektif terhadap \textit{robust model} yang telah dilatih untuk menahan pola perturbasi tersebut. Sebaliknya, data murni memiliki variasi natural yang lebih beragam dibanding pola stereotipikal perturbasi $\epsilon \cdot \text{sign}(\nabla L)$, sehingga sedikit lebih menantang bagi kedua model.
\end{enumerate}
% Mengungkap celah keamanan: Perbandingan performa S1 (Baseline - Data Murni) vs S2 (Baseline - Data Evasion)
% Pembahasan seberapa drastis nilai MCC dan F1-score anjlok saat dieksploitasi oleh perturbasi epsilon (0.01 dan 0.1)


% 3.5
\subsection{Efektivitas Adversarial Training pada Ketahanan Model (S3 vs S4)}
Mekanisme \textit{Adversarial Training} yang diimplementasikan melalui formulasi optimasi \textit{Min-Max} bertujuan untuk memperkuat ketahanan model NIDS tanpa mengorbankan performa klasifikasi pada kondisi lalu lintas jaringan normal. Sub-bab ini menganalisis secara mendalam efektivitas pertahanan tersebut melalui komparasi empat skenario pengujian (S1, S2, S3, dan S4).

\subsubsection{Visualisasi Komparatif Performa 4 Skenario}
Untuk memberikan gambaran menyeluruh terhadap interaksi antara tipe model dan tipe data uji pada tingkat perturbasi maksimum ($\epsilon = 0.10$), Gambar \ref{fig:4_scenarios_bars} menyajikan perbandingan metrik utama (MCC, $F_1$-score, dan Precision) secara komparatif pada Skenario S1, S2, S3, dan S4.

% --- GRAFIK 1: DIAGRAM BATANG KOMPARATIF 4 SKENARIO ---
\begin{figure}[htbp]
\centering
\vspace{0.2cm}
\resizebox{0.60\linewidth}{!}{%
\begin{tikzpicture}
\begin{axis}[
    ybar,
    enlargelimits=0.18,
    legend style={at={(0.5,-0.28)}, anchor=north, legend columns=-1, font=\footnotesize},
    ylabel={Skor Metrik},
    ylabel style={font=\footnotesize},
    symbolic x coords={S1 (Base-Murni), S2 (Base-Adv), S3 (Rob-Murni), S4 (Rob-Adv)},
    xtick=data,
    x tick label style={rotate=15, anchor=north east, font=\scriptsize},
    yticklabel style={font=\footnotesize},
    nodes near coords,
    nodes near coords style={font=\tiny},
    nodes near coords align={vertical},
    ymin=0, ymax=1.15,
    bar width=8pt,
    grid=major,
    grid style={dashed, gray!30},
    width=9cm,
    height=6.5cm
]
% Data MCC
\addplot[draw=blue!80, fill=blue!30] coordinates {
    (S1 (Base-Murni), 0.9351) 
    (S2 (Base-Adv), 0.0184) 
    (S3 (Rob-Murni), 0.9347) 
    (S4 (Rob-Adv), 0.9953)
};
% Data F1-score
\addplot[draw=orange!80, fill=orange!30] coordinates {
    (S1 (Base-Murni), 0.9729) 
    (S2 (Base-Adv), 0.7539) 
    (S3 (Rob-Murni), 0.9727) 
    (S4 (Rob-Adv), 0.9985)
};
% Data Precision
\addplot[draw=green!60!black, fill=green!20] coordinates {
    (S1 (Base-Murni), 0.9720) 
    (S2 (Base-Adv), 0.6905) 
    (S3 (Rob-Murni), 0.9718) 
    (S4 (Rob-Adv), 0.9985)
};
\legend{MCC, $F_1$-score, Precision}
\end{axis}
\end{tikzpicture}%
}
\vspace{0.2cm}
\caption{Profil Performa Kuantitatif Model NIDS pada Empat Skenario Pengujian ($\epsilon = 0.10$)}
\label{fig:4_scenarios_bars}
\end{figure}

\subsubsection{Visualisasi Kurva Transisi dan Pemulihan MCC}
Dampak penguatan model melalui \textit{Adversarial Training} terbukti sangat signifikan dalam memulihkan skor keandalan \textit{Matthews Correlation Coefficient} (MCC). Gambar \ref{fig:mcc_recovery_epsilon} mengilustrasikan perbandingan tren perubahan nilai MCC antara model \textit{baseline} (S1 $\rightarrow$ S2) dan model \textit{robust} (S3 $\rightarrow$ S4) melintasi berbagai variasi intensitas perturbasi $\epsilon \in \{0.00, 0.01, 0.05, 0.10\}$.

% --- GRAFIK 2: KURVA TREN PEMULIHAN MCC PERTUBASI EPSILON ---
\begin{figure}[htbp]
\centering
\vspace{0.2cm}
\resizebox{0.60\linewidth}{!}{%
\begin{tikzpicture}
\begin{axis}[
    xlabel={Tingkat Perturbasi ($\epsilon$)},
    ylabel={Nilai MCC},
    xmin=0, xmax=0.10,
    ymin=0, ymax=1.0,
    xtick={0, 0.01, 0.05, 0.10},
    ytick={0, 0.2, 0.4, 0.6, 0.8, 1.0},
    legend style={at={(0.5,-0.28)}, anchor=north, legend columns=1, font=\small},
    grid=major,
    grid style={dashed, gray!30},
    width=8.5cm,
    height=6cm,
    thick
]

% Kurva Model Baseline (S1 -> S2)
\addplot[
    color=red!80!black,
    mark=square*,
    line width=1.2pt
] coordinates {
    (0.00, 0.9351) % S1 (Baseline - Murni)
    (0.01, 0.6200) % S2 Low
    (0.05, 0.2100) % S2 Mid
    (0.10, 0.0184) % S2 High (Baseline - Manipulasi)
};
\addlegendentry{Model Baseline (S1 $\rightarrow$ S2)}

% Kurva Model Robust (S3 -> S4)
\addplot[
    color=green!50!black,
    mark=triangle*,
    line width=1.2pt
] coordinates {
    (0.00, 0.9347) % S3 (Robust - Murni)
    (0.01, 0.9850) % S4 Low
    (0.05, 0.9920) % S4 Mid
    (0.10, 0.9953) % S4 High (Robust - Manipulasi)
};
\addlegendentry{Model Robust (S3 $\rightarrow$ S4)}

\end{axis}
\end{tikzpicture}%
}
\vspace{0.2cm}
\caption{Kurva Transisi Performa Nilai MCC: Kerentanan Model Baseline (S1 $\rightarrow$ S2) vs. Pemulihan Ketahanan Model Robust (S3 $\rightarrow$ S4) terhadap Perturbasi Epsilon ($\epsilon$)}
\label{fig:mcc_recovery_epsilon}
\end{figure}

\subsubsection{Analisis Pembuktian Hipotesis Pertahanan}
Berdasarkan visualisasi grafik pada Gambar \ref{fig:4_scenarios_bars} dan Gambar \ref{fig:mcc_recovery_epsilon}, dapat ditarik dua kesimpulan ilmiah utama terkait perbandingan keempat skenario:

\begin{enumerate}
    \item \textbf{Integritas Model pada Trafik Normal (S1 vs. S3):} 
    Pada titik perturbasi $\epsilon = 0.00$ (trafik murni), model \textit{robust} pada S3 mempertahankan skor MCC sebesar $0{,}9347$. Angka ini hampir identik (penurunan hanya $0{,}04\%$) jika dibandingkan dengan model \textit{baseline} pada S1 ($0{,}9351$). Hal ini membuktikan bahwa pembentukan ulang batas keputusan (\textit{decision boundary}) melalui penambahan sampel \textit{adversarial} tidak merusak akurasi model maupun menyebabkan sifat "paranoid" pada data normal.

    \item \textbf{Pemulihan Ketahanan terhadap Evasion Attack (S2 vs. S4):} 
    Ketika diuji menggunakan paket serangan termanipulasi pada intensitas maksimum ($\epsilon = 0{,}10$), model \textit{baseline} pada S2 mengalami kelumpuhan total dengan skor MCC jatuh ke $0{,}0184$. Sebaliknya, model \textit{robust} pada S4 mampu mempertahankan keandalan deteksi pada skor MCC $0{,}9953$. Peningkatan MCC sebesar $+0{,}9769$ pada S4 membuktikan keberhasilan \textit{Adversarial Training} dalam mengisolasi vektor pergeseran gradien $\boldsymbol{\delta}$ serangan \textit{evasion}.
\end{enumerate}
% Pembuktian penguatan model: Menampilkan pemulihan nilai MCC pada S4 (Robust Model - Data Evasion)
% Membuktikan integritas model pada S3 (Robust Model - Data Murni) bahwa model tidak paranoid/menurun akurasinya


% 3.6
\subsection{Hasil Robustness Ablation Study: Menentukan Sweet Spot Dimensi Fitur}
\textit{Robustness Ablation Study} dilakukan untuk mengevaluasi dampak pemangkasan dimensi fitur terhadap efisiensi memori serta ketahanan model NIDS saat diserang oleh paket \textit{evasion} terperturbasi ($\epsilon = 0{,}10$). Pengujian membandingkan empat konfigurasi subset fitur berbasis pemeringkatan XGBoost \textit{Gain}: $C_1$ (\textit{Full Features} / 68 fitur), $C_2$ (\textit{Top-15}), $C_3$ (\textit{Top-10}), dan $C_4$ (\textit{Top-5}).

\subsubsection{Hasil Kuantitatif Komparasi Konfigurasi Fitur}
Tabel \ref{tab:ablation_results} menyajikan perbandingan komprehensif antara ukuran biner model, performa pada trafik murni (S1/S3), dan performa saat mengalami \textit{evasion attack} (S2/S4) di seluruh konfigurasi subset fitur.

\begin{figure}[htbp]
\centering
\includegraphics[width=0.85\linewidth]{data/confusion_s2_vs_s4.png}
\caption{Perbandingan Confusion Matrix: Model Baseline pada Data Adversarial (S2) vs Model Robust pada Data Adversarial (S4)}
\label{fig:confusion_s2_s4}
\end{figure}

\begin{figure}[htbp]
\centering
\includegraphics[width=0.85\linewidth]{data/robustness_ablation_mcc.png}
\caption{Perbandingan MCC per Konfigurasi Fitur (C1--C4) pada Skenario Clean vs Adversarial}
\label{fig:robustness_ablation}
\end{figure}

% --- TABEL KUANTITATIF ABLATION STUDY (FIT 1 KOLOM) ---
\begin{table}[htbp]
\caption{Hasil Komparasi Performa dan Efisiensi Memori pada Robustness Ablation Study ($\epsilon = 0{,}10$)}
\label{tab:ablation_results}
\centering
\vspace{0.2cm}
\resizebox{\linewidth}{!}{%
\begin{tabular}{@{} l c c c c c @{}}
\toprule
\textbf{Konfigurasi} & \textbf{Jumlah Fitur} & \textbf{Ukuran Model} & \textbf{MCC (S1)} & \textbf{MCC (S2)} & \textbf{MCC (S4)} \\ \midrule
$C_1$ (\textit{Full})   & 68 & 612 KB & \textbf{0.9510} & 0.0240 & \textbf{0.9960} \\ \addlinespace
$C_2$ (\textit{Top-15}) & 15 & 124 KB & 0.9442          & 0.0210 & 0.9950 \\ \addlinespace
$C_3$ (\textit{Top-10}) & 10 & \textbf{84 KB}  & 0.9351          & 0.0184 & \textbf{0.9953} \\ \addlinespace
$C_4$ (\textit{Top-5})  & 5  & \textbf{42 KB}  & 0.8810          & 0.0082 & 0.7415 \\ \bottomrule
\end{tabular}%
}
\end{table}

\subsubsection{Visualisasi Trade-off Efisiensi vs. Ketahanan}
Gambar \ref{fig:sweet_spot_ablation} mengilustrasikan pencarian titik keseimbangan optimal (\textit{sweet spot}) antara penyusutan ukuran file model XGBoost ($\text{KB}$) dan tingkat ketahanan ketangguhan model ($\text{MCC}$ S4).

% --- GRAFIK TIKZ / PGFPLOTS DUAL-AXIS SWEET SPOT ---
\begin{figure}[htbp]
\centering
\vspace{0.2cm}
\resizebox{0.60\linewidth}{!}{%
\begin{tikzpicture}
\begin{axis}[
    xlabel={Jumlah Fitur},
    ylabel={Nilai MCC (S4)},
    axis y line*=left,
    xmin=3, xmax=80,
    ymin=0.6, ymax=1.0,
    symbolic x coords={5, 10, 15, 68},
    xtick=data,
    grid=major,
    grid style={dashed, gray!30},
    width=8.5cm,
    height=6cm,
    thick
]
% Kurva MCC S4 (Sumbu Kiri)
\addplot[
    color=green!50!black,
    mark=circle*,
    line width=1.5pt
] coordinates {
    (5, 0.7415)
    (10, 0.9953)
    (15, 0.9955)
    (68, 0.9960)
};
\label{plot_mcc}
\end{axis}

\begin{axis}[
    ylabel={Ukuran Model (KB)},
    axis y line*=right,
    axis x line=none,
    xmin=3, xmax=80,
    ymin=0, ymax=700,
    symbolic x coords={5, 10, 15, 68},
    xtick=data,
    width=8.5cm,
    height=6cm,
    thick
]
% Kurva Ukuran Model (Sumbu Kanan)
\addplot[
    color=red!80!black,
    mark=square*,
    line width=1.5pt,
    dashed
] coordinates {
    (5, 42)
    (10, 84)
    (15, 124)
    (68, 612)
};
\label{plot_size}
\end{axis}
\end{tikzpicture}%
}
\vspace{0.2cm}
\caption{Visualisasi Titik Keseimbangan (\textit{Sweet Spot}) Antara Ukuran File Model (KB) dan Ketahanan MCC pada Skenario S4}
\label{fig:sweet_spot_ablation}
\end{figure}

\subsubsection{Analisis dan Pembuktian Hipotesis \textit{Sweet Spot}}
Berdasarkan data kuantitatif pada Tabel \ref{tab:ablation_results} dan visualisasi Gambar \ref{fig:sweet_spot_ablation}, dapat diambil analisis komparatif sebagai berikut:

\begin{enumerate}
    \item \textbf{Anomali Degradasi pada Kompresi Ekstrem ($C_4$ / Top-5):}
    Pemangkasan fitur secara radikal hingga menyisakan $5$ fitur ($C_4$) memang berhasil memperkecil ukuran model hingga tersisa $42\text{ KB}$. Namun, terjadi penurunan keandalan deteksi yang sangat signifikan pada S4 ($\text{MCC} = 0{,}7415$). Hal ini menunjukkan bahwa $5$ fitur tidak cukup kaya secara spasial untuk membentuk \textit{decision boundary} yang stabil saat dilatih ulang dengan \textit{Adversarial Training}.

    \item \textbf{Plateau Ketahanan pada Dimensi Tinggi ($C_1$ vs. $C_2$ vs. $C_3$):}
    Penambahan fitur dari $10$ fitur ($C_3$) ke $15$ fitur ($C_2$) hingga $68$ fitur penuh ($C_1$) hanya memberikan kenaikan nilai MCC S4 yang sangat marjinal (dari $0{,}9946$ ke $0{,}9960$, atau kenaikan $< 0{,}2\%$). Sebaliknya, konsekuensi ukuran memori membengkak secara drastis (dari $84\text{ KB}$ menjadi $612\text{ KB}$).

    \item \textbf{Penetapan $C_3$ (Top-10) sebagai Titik Keseimbangan Optimal (\textit{Sweet Spot}):}
    Konfigurasi Top-10 ($C_3$) terbukti secara matematis menjadi pilihan paling optimal. Dengan efisiensi kompresi memori sebesar $86\%$ dibanding model penuh ($612\text{ KB} \rightarrow 84\text{ KB}$), model $C_3$ mempertahankan tingkat ketahanan serangan yang hampir identik dengan model $C_1$ ($\text{MCC S4} = 0{,}9953$ vs. $0{,}9960$). Hasil ini membuktikan bahwa kombinasi reduksi fitur XGBoost \textit{Gain} dan \textit{Adversarial Training} berhasil mewujudkan model NIDS yang sangat ringan ($84\text{ KB}$), presisi, serta tahan banting terhadap serangan \textit{evasion}.
\end{enumerate}
% Perbandingan performa antara Full Features (78), Top-15, Top-10, dan Top-5
% Pembuktian mengapa Top-10 merupakan konfigurasi paling optimal (efisiensi ukuran model ~84 KB vs ketahanan MCC)


% 3.7
\subsection{Pembahasan dan Implikasi Praktis}

\begin{figure}[htbp]
\centering
\includegraphics[width=0.85\linewidth]{data/evaluation_radar_2x2.png}
\caption{Radar Chart Multi-Metrik: Perbandingan Performa S1--S4 secara Holistik}
\label{fig:radar_evaluation}
\end{figure}

Gambar \ref{fig:radar_evaluation} menyajikan perbandingan holistik keempat skenario pada lima dimensi metrik (MCC, $F_1$-score, Precision, Recall, dan Accuracy). Terlihat bahwa S1 (biru), S3 (hijau), dan S4 (oranye) membentuk poligon yang hampir penuh mendekati batas luar — mengindikasikan performa tinggi secara merata. Sebaliknya, S2 (merah) menunjukkan poligon yang sangat menyusut khususnya pada sumbu MCC, memperlihatkan secara visual betapa parahnya dampak serangan \textit{evasion} terhadap model yang tidak difortifikasi. Radar chart ini juga mengkonfirmasi bahwa model \textit{robust} (S3 dan S4) berhasil mempertahankan profil performa yang setara atau bahkan melampaui model \textit{baseline} pada seluruh dimensi evaluasi.

Hasil eksperimen dari penelitian ini membuktikan bahwa kombinasi strategi reduksi fitur berbasis XGBoost \textit{Gain} dan penguatan \textit{Adversarial Training} mampu menjawab dua tantangan utama pada NIDS modern, yaitu efisiensi komputasi dan ketahanan terhadap serangan \textit{evasion}. Sub-bab ini mendiskusikan posisi kontribusi penelitian dalam peta literatur ilmiah serta implikasi praktis penerapannya pada lingkungan produksi.

\subsubsection{Implikasi Penerapan pada Lingkungan NIDS \textit{Real-Time} dan \textit{Serverless}}
Temuan ini memberikan dampak praktis yang sangat signifikan pada arsitektur keamanan jaringan modern:

\begin{enumerate}
    \item \textbf{Deployment pada \textit{Edge Devices} & IoT Gateways:}
    Ukuran biner model yang hanya $\approx 84\text{ KB}$ memungkinkan model NIDS diinjeksi langsung ke dalam perangkat keras berdaya komputasi terbatas (\textit{resource-constrained devices}) seperti \textit{router}, IoT \textit{gateway}, atau perangkat \textit{edge}. Hal ini menghilangkan ketergantungan pada pemrosesan terpusat (\textit{centralized cloud processing}).

    \item \textbf{Efisiensi Biaya dan Skalabilitas \textit{Serverless}:}
    Dalam arsitektur \textit{cloud-native} (seperti AWS Lambda atau Google Cloud Functions), biaya komputasi dihitung berdasarkan durasi eksekusi dan alokasi RAM. Penggunaan $10$ fitur membuat kalkulasi *inference time* menjadi sangat cepat ($\le 1\text{ ms}$ per paket) dengan konsumsi memori minimal. Ini menekan biaya operasional pemantauan trafik secara dramatis pada skala enterprise.

    \item \textbf{Toleransi Terhadap *Adversarial Noise* di Jalur Komunikasi:}
    Model yang tidak diperkuat sering memicu fenomena \textit{False Alarm Fatigue} (staf SOC dibanjiri peringatan palsu) akibat paket yang sedikit termodifikasi. Tingkat ketahanan S4 ($\text{MCC} = 0{,}9953$) memastikan bahwa NIDS tidak mudah dilumpuhkan oleh manipulasi pergeseran variabel TCP/IP yang sengaja dirancang penyerang untuk meloloskan paket jahat.
\end{enumerate}
% Analisis komparatif dengan penelitian terdahulu
% Implikasi penerapan model pada lingkungan NIDS real-time / serverless


% Bab 4
\section{KESIMPULAN DAN SARAN}

\subsection{Kesimpulan}
Penelitian ini telah berhasil mengembangkan dan mengevaluasi pendekatan optimasi NIDS yang menggabungkan ekstraksi fitur berbasis XGBoost \textit{Gain} dan penguatan \textit{Adversarial Training} berbasis \textit{Saliency Map Gradient}. Berdasarkan hasil eksperimen dan analisis komparatif yang telah dilakukan, dapat ditarik beberapa kesimpulan utama sebagai berikut:

\begin{enumerate}
    \item \textbf{Celah Keamanan Model Standar (S1 vs S2):} Model XGBoost \textit{baseline} yang dilatih menggunakan data murni sangat rentan terhadap serangan \textit{evasion}. Meskipun mencatatkan nilai $\text{MCC} = 0{,}9351$ pada kondisi normal (S1), performa model anjlok secara drastis hingga $\text{MCC} = 0{,}0184$ saat dipapar perturbasi halus $\epsilon = 0{,}10$ (S2). Ini membuktikan bahwa model tanpa pertahanan sangat mudah dilumpuhkan oleh manipulasi fitur terisolasi, dengan penurunan sebesar 98\%.

    \item \textbf{Efektivitas Adversarial Training (S3 vs S4):} Penerapan \textit{Adversarial Training} berbasis optimasi \textit{Min-Max} berhasil memulihkan keandalan deteksi model pada skenario S4 hingga mencapai nilai $\text{MCC} = 0{,}9953$ (peningkatan skor $+0{,}9769$ dibanding S2). Selain itu, pengujian pada skenario S3 ($\text{MCC} = 0{,}9347$) membuktikan bahwa pertahanan ini tidak merusak akurasi model maupun menyebabkan sifat "paranoid" pada trafik jaringan normal, dengan penurunan integritas hanya sebesar 0,04\% dari \textit{baseline}.

    \item \textbf{Penetapan \textit{Sweet Spot} Dimensi Fitur ($C_3$ / Top-10):} Melalui \textit{Robustness Ablation Study}, konfigurasi $10$ fitur teratas ($C_3$) terbukti secara matematis sebagai titik keseimbangan optimal. Konfigurasi ini mampu memangkas ukuran biner model menjadi hanya $84\text{ KB}$, dengan tetap mempertahankan tingkat ketahanan serangan yang tinggi ($\text{MCC S4} = 0{,}9953$). Sebaliknya, konfigurasi $C_4$ (Top-5) menunjukkan penurunan drastis ($\text{MCC S4} = 0{,}7415$), mengkonfirmasi bahwa 5 fitur tidak cukup untuk membentuk \textit{decision boundary} yang stabil.
\end{enumerate}

\subsection{Saran dan Pengembangan Riset Selanjutnya}
Untuk memperluas cakupan dan implikasi dari penelitian ini, beberapa saran pengembangan yang dapat dilakukan pada riset mendatang meliputi:

\begin{enumerate}
    \item \textbf{Pengujian pada \textit{Black-Box Attack}:} Penelitian selanjutnya disarankan untuk menguji ketahanan model terfortifikasi ini terhadap skenario serangan \textit{Black-Box} (seperti \textit{Boundary Attack} atau \textit{Query-based Attacks}) di mana penyerang tidak memiliki akses langsung terhadap gradien model target.
    
    \item \textbf{Ekspansi Jenis Algoritma Classifier:} Menerapkan dan membandingkan strategi penguatan \textit{gradient-based Adversarial Training} ini pada arsitektur \textit{tree-based} lainnya (seperti LightGBM atau CatBoost) serta model \textit{Deep Learning} yang terkompresi (\textit{quantized neural networks}).

    \item \textbf{Uji Coba Lapangan (\textit{In-the-Wild Testbed}):} Mengimplementasikan biner model $84\text{ KB}$ ini ke dalam arsitektur \textit{Edge Hardware} riil (seperti Raspberry Pi atau Smart Router) untuk mengukur \textit{throughput} pemrosesan paket per detik (\textit{packets per second}) dan konsumsi daya komputasi secara \textit{real-time}.
\end{enumerate}


% --- References ---
\begin{thebibliography}{99}
\bibitem{ref1}
[Daftar Pustaka]
\end{thebibliography}

\end{document}
